\chapter{CASA}\label{casa}

\hfill\break

Il video della scatola nera in cui si vedeva l'asteroide distrutto dalla
bomba atomica era divenuto popolare tra l'equipaggio. L'equipaggio
sopravvissuto. Non era stata mia l'idea di far vedere alle persone il
video, non l'avevo neppure menzionato. Dovevano essere stati Desai o
Walorski a dire qualcosa, o qualcuno aveva chiesto a Skippy cosa fosse
successo e lui si era offerto di mostrarglielo. Osservare la distruzione
totale della base di ricerca sull'asteroide fu catartico per
l'equipaggio, come lo era stato per me. I medici dell'esercito ci dicono
che parlare degli eventi traumatici aiuta ad affrontare i sintomi del
disturbo da stress post-traumatico; tenere le cose dentro di te o
cercare di non ricordare cosa ti è successo, peggiora le cose. Skippy mi
disse che i sopravvissuti gli avevano chiesto se ci fosse un video
dell'assalto, girato dalle telecamere dei combot e delle tute kristang.
Un video c'era e io gli dissi di rilasciarlo a chiunque volesse
guardarlo. Giraud e io visionammo tutti i feed dei dati delle telecamere
e dei sensori, per renderci conto di cosa avevamo fatto bene e cosa
avevamo sbagliato e imparare la lezione.

Giraud mi rimproverò con moderazione di essere stato poco professionale,
far saltare in aria l'intero asteroide non sarebbe stato necessario per
la missione, l'avevo fatto perché ero arrabbiato, perché volevo che le
lucertole avvertissero il dolore che provavamo noi per avere perso così
tante persone. Così tanti umani. Forse l'avevo fatto mosso dalla
frustrazione e dal senso di colpa per non aver preso parte all'assalto.
Chissà. La psicologia non è la mia area di competenza.

L'assalto era stato un successo e un fallimento insieme. Dal punto di
vista delle perdite di vite umane, era stato un fallimento: avevo perso
più della metà del nostro equipaggio originario e alcuni dei feriti
avrebbero impiegato mesi per riprendersi, anche con l'incredibile
tecnologia medica thuranin. Dal punto di vista del piano di Skippy di
trovare un modo per localizzare il Collettivo, era stato un fallimento:
la cosa che avevamo riportato con noi era proprio quello che lui aveva
chiesto, non era danneggiata, non era inerte; Skippy riuscì ad
accenderla e leggerne i dati, solo che non gli disse nulla di utile. Non
sapeva se il difetto fosse nel dispositivo o dentro di sé, aveva la
brutta sensazione di essere in qualche modo bloccato rispetto a ricordi
che gli avrebbero permesso di accedere ai dati che riguardavano gli
Anziani. Lui stesso. Il luogo da cui proveniva, chi era. Non disse a
nessuno dei suoi dubbi e delle sue paure, solo a me. Da parte mia,
mantenni il suo segreto e cominciai a temere tacitamente che un giorno
Skippy potesse chiudersi o dare messaggio di errore, schermo blu, o
andare in standby, o qualsiasi cosa fosse successa alle antiche
intelligenze artificiali.

Era meglio che non si chiudesse con noi prima che fossimo tornati al
wormhole vicino alla Terra perché, in questo senso, la nostra missione
era stata un successo. Forse, sotto l'unico punto di vista che contava,
la nostra missione era stata un grande successo. Era stata un successo
tale che, anche se fosse sopravvissuta una sola persona per pilotare la
nave, l'assalto sarebbe valso il suo terribile costo. Avevamo rubato un
modulo di controllo del wormhole creato dagli Anziani, ed era intatto.
Skippy verificò che fosse del tutto funzionante, in grado di tenere fede
al suo nome e controllare un wormhole. E, cosa fondamentale, chiudere un
wormhole. Noi umani non avevamo idea di come funzionasse quella cazzo di
cosa, avremmo dovuto fare affidamento su Skippy per questo. Il punto era
che, grazie all'assalto, avevamo la capacità di chiudere quel wormhole.
Avevamo la capacità d'impedire ai kristang l'accesso alla Terra. I
miliardi di esseri umani là avrebbero ritenuto adeguato qualsiasi prezzo
la nostra ciurma di pirati avesse dovuto pagare.

Il fatto di doverci affidare a Skippy era un problema, era frustrato e
si aspettava che continuassimo la missione con l'equipaggio restante.
Quella era una conversazione che non volevo affrontare con gli altri in
ascolto: misi a tacere Skippy fino a quando il maggiore Simms non si
sarebbe seduta al comando, mentre Chang e Giraud erano entrambi in
infermeria. Dopo avere ripreso la Fiore a bordo, saltammo tre volte,
finché Skippy non fu certo che fossimo al sicuro, fuori dalla portata
della rete di sensori kristang. I motori di salto erano quasi scarichi,
capaci di appena un micro-salto di emergenza, ed eravamo sospesi in uno
spazio interstellare profondo, in attesa che si ricaricassero del tutto.

«Ti senti a tuo agio ora, Joe?», iniziò Skippy col suo tatto di sempre.
«Dobbiamo parlare della nostra prossima mossa. Riesco a vedere altre due
possibilità per contattare il Collettivo. Una si trova a soli tremila
anni luce di distanza, purtroppo arrivarci richiede un percorso di
wormhole traverso, ed è la meno probabile delle due opzioni. E poi, si
trova su un pianeta thuranin, sarebbe difficile entrare e uscire. La
seconda si trova a novemila anni luce di distanza, fuori dal territorio
dei thuranin, il che è un bene e un male, la specie che detiene quel
territorio è quella dei\ldots»

Cazzo. Non avevo neppure fatto in tempo a slacciarmi gli stivali. Era
meglio che dicessi solo quello che avevo in mente. «Skippy, prima che
continuiamo a cercare il Collettivo, andremo sulla Terra per rifornirci.
Non abbiamo abbastanza persone per continuare con la missione in questo
momento. Sei un genio, fai i conti.»

Skippy ebbe un'esitazione più breve che mai, di cui non mi sarei accorto
se non l'avessi conosciuto così bene. Aveva davvero fatto i conti?
Certo, millemila permutazioni. Forse un fantastilione. «Sei di cattivo
umore, Joe.»

«Sono un po' stanco. Un sacco di persone sono morte oggi, Skippy.» Mi
sfilai gli stivali scalciando, li avrei fatti volare attraverso il
compartimento, se la paratia non fosse stata a un solo metro di
distanza. «Esseri senzienti. Non dirmi cagate del tipo che siamo solo
scimmie e batteri, siamo senzienti. Siamo importanti.»

«Più di quanto immagini, Joe. Come mi hai chiesto, ho fatto due conti.
Il problema è che non sono uno stratega militare e ci sono un sacco di
variabili che non posso quantificare. Mi hai promesso che avremmo
trovato il Collettivo, insieme.»

«Skippy», mi sdraiai sul materasso come meglio potevo. «Io mantengo le
mie promesse. Spiegami come facciamo ad avere una ragionevole
probabilità di successo, con l'equipaggio che abbiamo ora, e ti
ascolterò. Io non la vedo. Ora abbiamo nove persone atte al
combattimento; io e due piloti dobbiamo restare a bordo dell'Olandese,
quindi rimangono sei elementi in azione, e due di loro sono
specializzati in logistica, non fanteria. Quell'asteroide era
l'obiettivo più debole tra le nostre opzioni, giusto?»

«È molto probabile, sì, i kristang non si rendevano conto di quello che
avevano, ecco perché gli oggetti di cui avevamo bisogno erano così poco
sorvegliati.»

«Poco sorvegliati?»

«Relativamente parlando.»

«Quella è la tua definizione di poco sorvegliato? Avevi il controllo
totale dei loro sensori e della maggior parte delle loro armi, eppure
siamo arrivati a uno scontro a fuoco. Anche i nostri combot
super-tecnologici sono stati massacrati. Uno degli altri due obiettivi
sarà più facile da conseguire?»

«Certo che no, no.»

«Pensi di non poterti fidare di noi, temi che, una volta arrivati sulla
Terra, non ce ne andremo mai più, che sarebbe la fine della missione.
Mettilo da parte per un minuto. Dimentica quello che ho detto. Il tuo
obiettivo è contattare il Collettivo. Usa il tuo potere di elaborazione
divino. Quali sono le probabilità di raggiungere il tuo obiettivo con le
risorse che abbiamo?»

Un'altra esitazione più breve che mai, forse più lunga stavolta? «Cazzo.
Hai ragione. Ero, forse, troppo sicuro di me, prima dell'assalto. Ora
che ho dati più estesi, ho rivisto la mia analisi. L'attuale probabilità
di successo della nostra missione è inferiore al 50\%. Per essere
precisi, è il 12\%. Questo è un livello di rischio inaccettabile.»

Ero così sollevato, che lasciai cadere la testa all'indietro e sbattei
la parte posteriore del cranio sul bordo di un armadietto. «Merda.»

«Merda nel senso che non puoi credere ai miei numeri, o merda nel senso
che credi ai miei numeri e sono peggio di quanto pensassi?»

«Merda, nel senso che ho sbattuto la testa contro questo stupido
armadietto un'altra volta.»

«Ah. Accidenti, hai voglia a calibrare il mio programma per leggere i
tuoi schemi e le tue espressioni vocali, continuo a sbagliare. Voi
esseri biologici siete esasperanti a volte.»

«Sì, come quando gli esseri biologici sviluppano la tecnologia e
costruiscono saccenti intelligenze artificiali.»

«Ti ho detto che non sono stato costruito da\ldots{} oh, lascia
perdere.»

Per abitudine consolidata, slacciai gli stivali e li misi di fronte a
me, vicino alla porta. In caso di emergenza, avrei potuto infilarci i
piedi e allacciarli molto in fretta, senza perdere tempo. Deformazione
militare, di nuovo. «Siamo d'accordo, allora. Torneremo sulla Terra,
giusto?» C'erano umani e rifornimenti umani su Paradiso, e Paradiso era
più vicino. È probabile che ci fossero anche un consistente contingente
di navi ruhar in orbita, soldati ruhar a terra e un gruppo di battaglia
jeraptha in agguato alle estremità del sistema. Paradiso non era un
posto in cui volevo tornare.

«Siamo d'accordo, cazzo. Ah, ho aspettato un milione di anni per
contattare il Collettivo, un breve ritardo non è un grosso problema,
credo.» La sua voce sembrava tutt'altro che convinta. «Bleah. Questo
significa che andremo a visitare quella palla di fango infestata da
scimmie che chiami casa.»

«Casa dolce casa, Skippy. Togliamoci il pensiero, così non ne
discuteremo più tardi. Come possiamo fare in modo che tu ti fidi del
fatto che noi, una volta arrivati sulla Terra, non ci resteremo? Con
``noi'' non intendo me, io manterrò la promessa che ti ho fatto.» Non
era una domanda retorica, io stesso ero preoccupato. Avere una nave
madre thuranin e una fregata kristang in orbita sarebbe stato molto
allettante per i governi della Terra, sarebbero stati tentati di tenerci
lì. E avrebbero addotto argomentazioni molto valide per mantenere le
nostre due astronavi in orbita: per esaminare la loro tecnologia, per
proteggere la Terra nel caso in cui qualche nave kristang rimasta
indietro fosse ancora in giro, perché rimandare le nostre navi a vagare
per la galassia, finché Skippy non avesse trovato un Collettivo che
poteva esistere oppure no, era pura idiozia. Non è che volessi tornare a
vagare tra le stelle con Skippy, per niente. Quello che volevo era che
la vita tornasse alla normalità, era togliermi le stupide aquile
d'argento, tornare a essere un soldato semplice, finire il mio periodo
di servizio, tornare a casa e vivere una normale vita umana. E le
bambine vogliono che Babbo Natale porti loro un pony. Dovevamo entrambi
accettare la realtà. Non c'era più una vita ``normale'' per l'umanità,
certo non per me. Se io o Skippy non fossimo riusciti a trovare un modo
per costringere i governi della Terra ad autorizzare l'invio
dell'Olandese Volantein orbita, non avrei mai potuto mantenere la mia
promessa.

«Ah, quello. Nessun problema, Joe. Dopo che avremo attraversato l'ultimo
wormhole che porta alla Terra, interromperò in via temporanea la sua
connessione alla rete. Questo lo disabiliterà finché non si sarà
resettato, il che richiederà abbastanza tempo perché noi arriviamo sulla
Terra, carichiamo volontari e rifornimenti e torniamo al wormhole.»

«Uhm. È una grande idea, Skippy.» Perché non ci avevo pensato? Skippy mi
aveva detto che poteva interrompere in via temporanea un wormhole da
solo, il modulo di controllo era necessario per disattivarne uno in modo
permanente. «Avremo tempi stretti, allora?»

«Stretti, sì. Una volta che avrò interrotto quel wormhole, i thuranin
avranno una gran voglia di sapere cosa diavolo sia successo. Non appena
si sarà resettato, delle navi lo attraverseranno. Dato che i thuranin
sono dei piccoli, sospettosi, paranoici figli di puttana verdi, posso
garantire che alcune delle loro navi faranno visita alla Terra, e non
saranno lì per usare i loro buoni omaggio da Starbucks. I thuranin hanno
un protocollo stabilito per l'investigazione di una nuova connessione
wormhole: ci inviano attraverso una forza pesante di navi da battaglia,
con incrociatori di scorta. Non vogliamo scontrarci con un carro da
battaglia thuranin, una di quelle navi potrebbe polverizzare l'Olandese
in particelle subatomiche, e persino io potrei solo rallentarli.» Skippy
aveva detto più volte che, pur con tutta la sua apparente potenza, una
nave madre non era in primo luogo una nave da guerra, ma un mezzo di
trasporto a lungo raggio per navi da guerra. Durante un combattimento,
le navi madri fuggivano e lasciavano combattere le loro scorte.

«Perfetto, Skippy. Per favore, traccia una nuova rotta di salto per la
Terra.» Saremmo rimasti sulla Terra solo il tempo di un caffè o, per
quanto mi riguardava, un cheeseburger. Ne sentivo già il sapore.
L'antica intelligenza artificiale degli Anziani mi aveva sorpreso di
nuovo, mi ero preparato a una lunga e accesa discussione per decidere se
avremmo continuato a cercare il Collettivo, o saremmo andati a
rifornirci sulla Terra. Skippy non reagiva alle discussioni, reagiva ai
fatti e alla logica. «Ehi», il mio dito si posò sull'interruttore per
spegnere la luce, «per curiosità, secondo i tuoi calcoli, quali erano le
probabilità di successo dell'assalto appena concluso?»

«In origine, 37,6\%. Poi, dopo che l'equipaggio ha dimostrato competenza
con i combot e il capitano Giraud ha sviluppato il suo piano d'assalto,
ho calcolato le probabilità al 51,1\%.»

«Appena il 50\%? Non pensavi che valesse la pena menzionarlo?!»

«Sopra il 50\%. Sopra è sopra. Non l'hai chiesto. E il tuo fascicolo
indica che la matematica non è il tuo forte.»

Discutere del passato non era produttivo. «I motori di salto saranno
carichi del tutto tra due ore e mezza, giusto? Svegliami tra due ore.»

«Incredibile. Sono la sensibilità più avanzata della galassia e mi stai
usando come sveglia.»

«Due ore, Skippy. Buonanotte.»

\hfill\break

I componenti dell'equipaggio furono sollevati ed entusiasti nel sentire
che ci saremmo fermati sulla Terra prima di continuare la missione.
Almeno, ne furono felici per un paio di giorni. Poi, durante uno dei
miei turni di servizio, Simms fece un cenno per attirare la mia
attenzione, io annuii e le feci segno di venire sul ponte di comando. In
quel momento, non stava succedendo niente di importante, eravamo a un
anno luce da una nana rossa di poco conto, a tre anni luce dal wormhole
più vicino, in attesa che i motori si ricaricassero. «Colonnello», disse
il maggiore Simms, «ora che stiamo per chiudere il wormhole, tra
l'equipaggio girano dei ripensamenti.»

«Cosa?», esclamò Desai dalla poltrona del pilota, e anche Walorski si
voltò.

Io ebbi la stessa reazione: «Maggiore, il fine dell'intera missione è
bloccare alle lucertole l'accesso alla Terra».

«Sì, signore, non è questo il problema. La gente teme che la
Expeditionary Force su Paradiso non potrà mai più tornare sulla Terra,
dopo la chiusura del wormhole; non riceveranno spedizioni di cibo o
forniture mediche. Li abbandoneremo. Per sempre.» Le sue mani si
strinsero a pugno, mostrando la sua ansia. «Questo non sta bene
all'equipaggio. O a me. Signore.»

Questa era una conversazione che sapevo imminente e che non ero
impaziente di affrontare. L'argomento mi era rimasto nel sottoscala
della mente fin da quando Skippy mi aveva parlato dell'idea di chiudere
il wormhole. Prima che potessi trovare una risposta, Skippy prese la
parola: «Non preoccuparti, maggiore Tammy. La vostra Expeditionary Force
è già stata abbandonata, e non possiamo farci niente, quindi non siamo
responsabili. Ho intercettato dei messaggi in cui i thuranin
comunicavano al clan Vento Bianco che non avrebbero sostenuto ulteriori
tentativi di riconquista di Paradiso, il pianeta non vale lo sforzo e le
forze thuranin sono impegnate altrove nel settore. I kristang avevano
smesso di portare rifornimenti dalla Terra ancor prima che i ruhar si
riprendessero il pianeta. I kristang non faranno alcuno sforzo per
evacuare gli umani da Paradiso e i ruhar non hanno accesso alla Terra,
né capacità di trasporto in questo settore. Sono da soli».

Mi aveva fatto incazzare di nuovo: «Cazzo, Skippy, non devi sembrare
così allegro!».

«I fatti non sono allegri o cupi, colonnello Joe», Skippy aveva un tono
difensivo nella voce, che mi sorprese. «Sono solo fatti. Il fatto è che,
che chiudiamo il wormhole o no, l'Unef è bloccata su Paradiso per il
prossimo futuro. I ruhar forniranno integratori alimentari fino a quando
non si potrà fare la raccolta nei campi da voi coltivati. Penso.»

«Skippy, non sei d'aiuto. Maggiore, non ho ignorato l'Unef, non volevo
parlarne, perché Skippy ha ragione al 1.000\% su questo. Non possiamo
fare nulla per aiutare l'Unef. Possiamo aiutare i miliardi di umani
sulla Terra. E se qualcuno ha la brillante idea che dovremmo consegnare
Skippy ai ruhar, perché i criceti trasportino in cambio la gente
dell'Unef sulla Terra, scordatevelo.»

Skippy fu felice di sentirlo: «Grazie, Joe, apprezzo».

«La programmazione di Skippy lo metterebbe in standby in presenza di
specie interstellari evolute, perciò, se lo dessimo ai criceti,
perderemmo la capacità di chiudere il wormhole. O qualsiasi altra cosa.»

«Ah», Skippy sembrava davvero sorpreso. «Che peccato, colonnello Joe,
per un attimo ho pensato che stesse esprimendo un tantino di lealtà nei
miei confronti.»

Sospirai prima di riuscire a fermarmi: «Skippy, hai detto chiaro e
tondo, ogni volta che hai potuto, che sei un potente super-essere e che
noi per te siamo batteri. Questa non è un'amicizia tra te e me, è
un'alleanza tra specie, culture o qualunque cosa sia. Noi ti siamo utili
e tu ci sei utile. Quando io faccio un accordo con te, gli mantengo
assolutamente fede. Quello che non so è se tu ritieni che valga la pena
mantenere fede a un qualsiasi accordo fatto con dei batteri».

«Uhm.»

«Sì, uhm.» Con la coda dell'occhio, percepii un avvertimento da parte di
Simms, con tutta probabilità temeva che rischiassi di far incazzare
Skippy. Non conosceva quella testa di cazzo come la conoscevo io.

«Mi sembra giusto, colonnello Joe. Voi non lo sapete dato che siete,
dopotutto, batteri, ma in realtà sono affidabile al 100\% quando faccio
una promessa. Possiamo lavorarci su.» Sembrava ancora ferito. Mi
chiedevo quale parte della sua capacità fosse dedicata a una sub-routine
di emulazione delle emozioni. Era quasi convincente.

«Abbiamo finito, maggiore?», chiesi a Simms. «Non c'è niente che
possiamo fare per aiutare la Expeditionary Force su Paradiso. L'Unef è
venuta qui per proteggere la Terra; quella missione si è incasinata
perché le lucertole ci hanno mentito, ma noi possiamo ancora proteggere
il nostro pianeta. È questa la nostra missione. Se mai ci fosse un modo
per ristabilire il contatto con l'Unef, per me andrebbe benissimo, ma va
oltre le nostre capacità.»

Simms annuì con una smorfia: «Vuole che parli all'equipaggio?».

Questo avrebbe significato prendere la via d'uscita del codardo. «No, me
ne occuperò io.»

Fui maggiormente diplomatico quando parlai all'equipaggio più tardi quel
giorno; spiegai la situazione, ascoltai con simpatia, dissi che avevo
anch'io amici su Paradiso. Quello che non feci fu farmi forte del mio
rango e dichiarare che ero il comandante e avremmo fatto quello che io
pensavo fosse meglio. Per lo più ascoltai e lasciai che la gente
parlasse. Tutti capivano che non c'era nulla che in pratica potessimo
fare per gli umani su Paradiso, e quello che davvero li disturbava era
il senso di colpa. Il senso di colpa perché saremmo, si sperava, andati
a casa e avremmo goduto di qualsiasi comodità terrestre vi si trovasse
ancora. A casa con amici e familiari e cheeseburger. Mentre la gente
dell'Unef sarebbe rimasta a sperare che i ruhar decidessero di deviare
dal loro sforzo di guerra risorse sufficienti a tenerla in vita. Anche
se i ruhar avessero deciso di nutrire i loro ex nemici, gli umani
tecnologicamente arretrati, i kristang avrebbero potuto disturbare le
loro spedizioni abbastanza da impedire ai rifornimenti di arrivare su
Paradiso, e i ruhar avrebbero potuto non essere in grado di provvedere
all'Unef.

Alla fine, ricordai alla nostra non tanto allegra banda di pirati che
saremmo tornati a casa in una situazione sconosciuta, che non sapevamo
quanti kristang e navi si trovassero sul lato della Terra rivolto verso
il wormhole, che la Terra avrebbe potuto non essere più il paradiso blu
e verde che ricordavamo, che avremmo potuto essere costretti ad aprirci
la strada combattendo contro forze nemiche superiori per numero. E che,
una volta chiuso il wormhole, quando i kristang sulla Terra si sarebbero
resi conto di non poter tornare a casa, avrebbero potuto essere tentati
di dimenticare Le Regole e usare armi proibite contro la popolazione
umana. Su questo, gli occhi della gente si strinsero e le mascelle si
sporsero. «La nostra missione non sarà finita quando avremo chiuso il
wormhole dietro di noi. Quello sarà l'inizio.»

«Colonnello», chiese il soldato Putri. Il Putri americano, non quello
indiano. «Qual è il piano nel caso in cui sulla Terra ci fosse una forza
considerevole di navi kristang?»

«In quel caso combatteremo. Non servirebbe chiudere il wormhole se le
lucertole sulla Terra potessero ancora distruggere la nostra casa.
Combatteremo come possiamo, finché i kristang non saranno più una
minaccia per la Terra. L'Olandese non è una nave da battaglia, ma
abbiamo armi e capacità di salto superiori. Se le lucertole vogliono
combattere, facciamo loro sputare sangue. Combatteremo finché non
saranno distrutte, o fino al nostro ultimo respiro. Questo è il mio
piano.»

\hfill\break

Un paio d'ore dopo, mi stavo togliendo le scarpe nel mio minuscolo
dormitorio, quando Skippy parlò dall'altoparlante sul soffitto.

«Colonnello Joe, dobbiamo parlare.''

«Oh», gemetti, «non parliamo già abbastanza? Puoi aspettare?»

«No e poi no. Siamo vicini al passaggio del wormhole per la Terra, ho
bisogno che tu stia a sentire una cosa prima di programmare il lancio.»

Ehm\ldots{} oh. Il suo tono mi fece rizzare le antenne. Problemi in
vista. Che cosa mi aveva taciuto quella piccola, lucida testa di cazzo
stavolta? Ritornando coi piedi sul pavimento, mi strofinai la faccia per
simulare uno stato di semivigilanza: «Sono tutt'orecchi. Cosa c'è?».

«È stato commovente il discorso che hai tenuto, sul fatto di dover
combattere i kristang anima e corpo, fino all'ultimo respiro.»

«Mi tieni sveglio per farmi i complimenti per il mio discorso?»

«No, a essere sincero, come discorso era al massimo di terzo livello.
Del tutto derivativo. Stai parlando della possibilità che si entri nella
prima battaglia spaziale dell'umanità e il meglio che sai dire sono
degli stupidi cliché? Potevi almeno aggiungere qualche citazione
esagerata da Patton, o qualcosa del genere.»

«Cristo santo, Sk\ldots»

«Quello che voglio dire è che potrai anche avere intenzione di
combattere anima e corpo, ma una cosa contro cui non dovrai combattere è
questa nave. Ho bisogno dell'Olandese per contattare il Collettivo. Non
permetterò che venga messa a rischio. E, dato che ho bisogno di un
umano, di un equipaggio umano vivo che la piloti per me, non ti
permetterò di mandare tutti a combattere. Specie se è probabile che tu
abbia la peggio.»

«Merda», risposi a denti stretti. «Che vuol dire che non me lo
permetti?»

«Se vuoi fare qualcosa di stupido che metterà in pericolo l'Olandese o
mi priverà di un equipaggio, non collaborerò alla gestione della nave.
Questo significa che non programmerò salti, né caricherò le rotte nel
pilota automatico, né preparerò e punterò le armi. Il capitano Desai ha
imparato solo a fare minime manovre nello spazio normale e posso
bloccare anche quei comandi. Passando attraverso il wormhole e
chiudendolo dietro di noi, correrò il rischio che tu decida di non
aiutarmi a trovare il Collettivo dopo che avremo fatto rifornimento
sulla Terra. Abbiamo un accordo, Joe, mi aspetto che tu lo onori.»

«Cazzo.»

«Che parola multiuso, Joe.»

«Intendevo della serie: cazzo, sì, non ho dimenticato il nostro accordo,
cui manterrò assolutamente fede. Hai qualche idea geniale su come
rifornirci, se i kristang hanno una task force di navi nell'orbita
terrestre?»

«Questo scenario è improbabile, dato il prezzo che i thuranin stanno
facendo pagare al clan Vento Bianco per il trasporto da e verso la
Terra. Mi aspetto che i kristang abbiano inviato solo una guarnigione di
una minima forza, il loro vantaggio tecnologico è così schiacciante che
non avrebbero bisogno di molte truppe per tenere sotto controllo il
vostro arretrato pianeta. Dalle comunicazioni dei thuranin che ho
intercettato ho saputo che importanti elementi della flotta si stanno
preparando per una battaglia attorno a un gruppo di wormhole all'altra
estremità di questo settore. La loro attenzione è concentrata altrove.»

«Il punto è, Skippy, che non abbiamo idea di cosa troveremo quando
arriveremo sulla Terra. Se scopriamo che le lucertole stanno
terrorizzando il nostro pianeta, non puoi aspettarti che questo
equipaggio scappi e si metta a vagare per la galassia con te. Ci serve
un'altra opzione.»

«Sono aperto ai suggerimenti, Joe. Sei tu il soldato.»

Come avrei potuto escogitare una strategia di battaglia spaziale, con la
mente stanca e nessuna esperienza o addestramento in quel genere di
combattimento? O in qualsiasi tipo di combattimento nave per nave? Ah,
sì, perché indossavo aquile d'argento, ecco come. «Va bene, che ne dici
della Fiore? Quella nave non ti serve, vero? Ti va bene se mando una
parte dell'equipaggio sulla Fiore a combattere i kristang? Ti prometto,
Skippy, che non ti lascerò qui. Rimarrò a bordo dell'Olandese.»

«La Fiore è in qualche modo utile, perché avere attraccata una nave
kristang, in particolare una nave kristang danneggiata in battaglia, è
uno stratagemma efficace. Ma non è essenziale. Ok, puoi staccare la
Fiore, io programmerò anche un salto per te. Non so che cosa ti aspetti
di buono dal tuo equipaggio non addestrato, con una nave che non sa
guidare.»

«Allora ho bisogno che inizino l'addestramento, subito». Avevamo anche
bisogno di una strategia. O, no? Presi in considerazione quello che
aveva detto Giraud sui piani di battaglia flessibili e aveva senso in
tutto e per tutto. Non avremmo potuto fare piani fino a quando non
avessimo avuto informazioni. Qualsiasi informazione. «Che ne dici di
questo, Skippy? Possiamo far saltare l'Olandese abbastanza vicino da
vedere quali forze i kristang hanno disposto attorno alla Terra, ma
abbastanza lontano da essere al sicuro? A quel punto, potremmo staccare
la Fiore, oppure, se concordi che il rischio è minimo, far saltare in
orbita l'Olandese.»

«Mmh. Suona sospettosamente ambiguo. Rischio minimo non equivale a
nessun rischio. Ah, che cazzo, perché no? Sono già annoiato. Solo per
te, Joe, concordo su questo: possiamo saltare direttamente in orbita ed
esaminare la situazione da lì e saltarcene fuori di nuovo se giudico
troppo grande il rischio. Non che non mi fidi di te o del tuo
equipaggio, ma programmerò il pilota automatico con un timer, per farci
saltare di nuovo fuori a meno che io non annulli l'ordine.»

«Affare fatto», mi affrettai ad accettare prima che Skippy potesse
cambiare idea. Ogni rischio che la mia lattina di birra
super-intelligente avesse giudicato troppo grande avrebbe significato
che era meglio per noi ritirarci e riprendere comunque in considerazione
le opzioni a disposizione. Forse avremmo potuto lanciare un paio di
proiettili ai kristang prima di partire, come sveglia. Una nave madre
thuranin che saltava in orbita sulla Terra e faceva esplodere un paio di
navi kristang senza preavviso avrebbe gettato le lucertole nel panico.

«Per la tua formazione in tattiche di combattimento spaziale, se i
kristang sono sulla Terra in forze, è meglio saltare in orbita piuttosto
che saltare fuori di un paio di unità astronomiche. Persino io non sono
in grado di nascondere l'esplosione di raggi gamma quando usciamo da un
salto. Dovremmo aspettare lì e avviare scansioni di sensori ad ampio
raggio, e i kristang sulla Terra sarebbero allertati del nostro arrivo.
Saltare loro addosso ci permetterà di coglierli alla sprovvista e magari
pompargli un paio di missili su per il culo prima che possano reagire.
Allora potremo saltare via, se ne abbiamo bisogno.»

«Il combattimento spaziale sembra complicato.» Ripensai a quando avevo
ascoltato il pilota di Pollo parlare di combattimento aereo dopo la
nostra prima simulazione di guerra, su Campo Alfa.

«E certo! E poi c'è il fattore Skippy.»

Una parte di me voleva evitare di abboccare: «Il fattore Skippy?».

«Sai, la mia incredibile grandiosità.»

«Ah, certo.»

«Non ne sei convinto? Nello specifico, mi riferisco alla mia capacità di
prendere il controllo dei sistemi kristang da remoto, usando il
nano-virus thuranin incorporato nei loro sistemi. Per farlo, devo essere
a circa un secondo luce dalla nave nemica.»

«Un secondo luce? Questo è, ehm\ldots», la luce viaggia, ehm\ldots{}
uhm\ldots{} stavo cercando di immaginare me stesso in una classe.

«Lascia che ponga fine alle tue sofferenze, prima che ti vada a fuoco il
cervello, scimmiotto. Un secondo luce è all'incirca la distanza tra la
Terra e la vostra luna.»

«Ah», pensai un momento, «ero convinto che le navi non dovessero saltare
così vicino a un pianeta, non è così?»

«La maggior parte delle navi non può, il pozzo gravitazionale di un
pianeta distorce il campo di salto all'ingresso, rendendo la navigazione
di salto imprevedibile, e la distorsione del campo può danneggiare i
motori di salto, anche fare a pezzi una nave. Ma una nave dove ci sono
io che controllo i motori di salto compensa la distorsione di campo.
Voilà, il fattore Skippy.»

«Impressionante», dovetti ammettere.

«Come? Niente battute sarcastiche?»

«No, sarai anche uno stronzo arrogante, ma la tua grandiosità è
legittima.»

«Non è vantarsi se è vero.»

«Sì, questa l'ho già sentita. Ehi, è per questo che ci hai fatto saltare
così vicino a quel gigante gassoso?»

Ci fu un'esitazione più breve che mai: «Può darsi che fosse troppo
vicino. Non avevo finito di mettere a punto i merdosi motori di salto di
questa nave. Non accadrà più. Ehi, bella chiacchierata. Devi dormire un
po', possiamo parlarne dopo, eh?».

Mi sdraiai sul letto, cercando di nascondere un sorriso. La scimmia
aveva messo a disagio l'intelligenza artificiale. Dovevo ricordarmelo.

La transizione attraverso l'ultimo wormhole andò liscia, ogni mio timore
di poterci imbattere in una nave thuranin in partenza era superfluo.
Skippy non rilevò nessuna nave nella zona. Quando fummo al sicuro fuori
dal wormhole, Skippy ne interruppe la connessione alla rete e lo spense.
Per fortuna, aveva caricato una nuova applicazione in evidenza sulla
schermata iniziale di tutti i nostri zPhone, era un semplice orologio.
Un conto alla rovescia fino al reset del wormhole. Tutti ricevemmo il
messaggio, forte e chiaro.

\hfill\break

Nel mio alloggio, stavo cercando di attaccare i gradi da sergente sulla
giacca di una delle mie uniformi, Skippy aveva fabbricato i galloni per
me. Una volta raggiunta la Terra, il mio rango teatrale di colonnello
sarebbe stato annullato e io sarei tornato al mio grado regolare di
sergente dell'esercito. Una parte di me, una gran parte di me, temeva
quello che gli alti ufficiali dell'esercito avrebbero pensato del fatto
che avevo indossato le aquile d'argento, e di tutti i casini che avevo
combinato, o delle decisioni discutibili che avevo preso. Non appena la
Terra sarebbe comparsa nello schermo, quelle aquile da guerra d'argento
sarebbero entrate in una scatola, e non ne avrei fatto una tragedia.
Chang, che era stato rilasciato dall'infermeria ora che le sue costole
rotte erano in via di guarigione, accettò la mia proposta di prendere il
comando quando avremmo raggiunto la Terra, con l'eccezione che io
restavo capitano della nave. Secondo me, Chang non era ancora a suo agio
con Skippy. Comunque non aveva molta importanza, una volta che avremmo
preso contatto con le autorità sulla Terra, il nostro equipaggio pirata
avrebbe risposto al comando del pianeta. Sempre che non fossimo stati
costretti ad aprirci la strada combattendo contro una flotta di navi da
guerra kristang. Questo pensiero m'innervosiva.

C'era un'altra cosa che m'innervosiva, aspettai di poterne parlare con
Skippy in privato. Una pausa dagli sprint di corsa lungo la chiglia
dell'Olandese era una buona occasione. «Senti, grande e potente Oz, mi
dispiace davvero di averti chiamato Skippy. Mi sento un idiota ora, non
sapevo quanto tu fossi potente e non volevo mancarti di rispetto. Col
governo faremmo una figura di merda se venissero a sapere che ti ho
chiamato Skippy, perciò come potrei chiamarti? Signore Dio Onnipotente è
ancora valido, nel caso ci stessi pensando.»

«Mi va bene Skippy. Mi piace.»

«Sei sicuro?» Non capivo se stesse scherzando.

«Sì. Skippy è un soprannome, giusto?»

«Credo di sì.» Non conoscevo nessuno che si chiamasse Skippy.

«E i soprannomi possono essere termini derisori, il che, ammettilo, non
è possibile quando siete voi forme di vita inferiori a parlare di me.»

«Certo che no», strabuzzai gli occhi.

«Vorrei vedere! Ma un soprannome può essere anche indicativo di
accettazione, di appartenenza, di adesione al gruppo dei ragazzi fighi.»

Ragazzi fighi? Un essere d'incredibile potenza, vecchio un milione di
anni, voleva essere uno dei ``ragazzi fighi''? «Certo, vada per Skippy.»

«Le persone che mi chiamano Skippy ricorderanno di continuo a voi
scimmie che non somiglio per niente a nessun altro picchiatello che
abbia mai vissuto su quel miserabile schifoso pianeta e la cosa vi farà
notare la vostra totale trascurabilità rispetto a me, molto meglio di
qualsiasi nome destinato a evocare rispetto. E, sul serio, pensi davvero
che voi imbranati sareste davvero capaci di darmi il rispetto che
merito? Il nome Skippy è appropriato; è stata una reazione difensiva da
parte tua, a qualcosa che va ben oltre la tua comprensione.»

«È stata una reazione offensiva al fatto che sei uno stronzo!»

«O quello. Come vuoi.»

Prima di fare il salto finale nel sistema solare di casa nostra, ordinai
una smobilitazione, per assicurarmi che tutti fossero ben riposati e
tutti i nostri sistemi e le attrezzature fossero resettati e pronti al
rock`n'roll. In particolare, ero preoccupato per le armi dell'Olandese,
per quanto relativamente deboli.

\hfill\break

Chang era in piedi, e più in fretta di quanto mi aspettassi, anche con i
trattamenti miracolosi thuranin: il ragazzo aveva le costole rotte e un
polmone in parte perforato. Entrò nel Cic, dove c'era anche Simms,
mentre io ero sulla poltrona di comando. «Colonnello Chang, non dovrebbe
essere in infermeria?», chiesi. Skippy non aveva detto nulla sul fatto
che avrebbe rilasciato Chang dalla terapia.

Chang si tolse la maglietta, aveva una plastica nera attorno alle
costole. «Mi stanno curando, Skippy ha detto che andare in giro
aiuterebbe i tessuti a adattarsi mentre guariscono.»

«Va bene, basta che non pensi di tornare in servizio. Vacci piano per un
paio di giorni, d'accordo?»

Chang trasalì, ancora sofferente quando si muoveva: «D'accordo.
Colonnello, ho sentito che è stata sua l'idea di usare le nostre radio
per triangolare la posizione del nodo di comunicazione. È stata una
trovata eccellente, forse quello che ci ha salvati tutti. Se avessimo
dovuto smistare tutta quella spazzatura da soli, saremmo rimasti
intrappolati».

Ricevere l'ammirazione di Chang faceva molto bene, peccato che fosse
malriposta. Purtroppo, dovetti spiegare cos'era successo davvero:
«Grazie, non era niente di brillante, era ovvio. Tutto quello che ho
fatto è stato chiedere se il nodo di comunicazione potesse trasmettere,
e mi pareva che dovesse poterlo fare, se il suo scopo sono, sai, le
comunicazioni. Skippy è incredibilmente intelligente, ma è anche
distratto e non pensa a cose che noi riterremmo ovvie. Per esempio, non
ci ha detto dei combot finché non gli ho chiesto come combattessero i
thuranin. Tienilo sempre a mente quando hai a che fare con Skippy, è
solo che non pensa al nostro livello».

\hfill\break

Saltammo nell'orbita terrestre, Skippy riferì che c'erano solo due navi
kristang, una fregata e una nave da trasporto truppe che era anche la
loro ammiraglia. Regolai lo schermo per zoomare sulla grande nave da
trasporto truppe kristang che riempì uno dei visori. Era maledettamente
grande, anche se continuavo a dimenticare quanto più grande fosse
l'Olandese, che avrebbe potuto trasportare dozzine di quelle navi
kristang attraverso anni luce.

«Avrò bisogno di stabilire una connessione con il nano-virus
sull'ammiraglia kristang, craccare la loro crittografia multilivello,
prendere il controllo dei loro computer e bloccare loro l'accesso.»

«Quanto tempo ci vorrà?»

«L'ho fatto mentre parlavo, nell'intervallo tra ``stabilire'' e
``connessione''.»

«A nessuno piacciono le ostentazioni, Skippy.»

«Posso andare più piano, se vuoi, ma è probabile che mi annoierei e
perderei di vista quello che avrei dovuto fare dopo un paio di
picosecondi.»

«Non possiamo permettercelo. Quali risorse hanno sulla Terra?»

«Solo quelle due navi: la nave da trasporto truppe laggiù, che è la loro
ammiraglia, e una fregata in orbita polare, al momento sopra Sumatra,
dall'altra parte del pianeta. La nave da trasporto ha ventiquattro
navicelle d'assalto di vario tipo, di cui tre sono a bordo ora, una è in
orbita in avvicinamento e le altre sono sparse per tutto il pianeta.
Anche la fregata ha a bordo due piccole navicelle. Ci sono insediamenti
difensivi dei thuranin in cima agli ascensori. E i kristang hanno una
costellazione di diciassette satelliti maser in orbita per gli attacchi
a terra.»

«Ok», feci un lungo respiro. «Siamo abbastanza sicuri da restare qui?
Annulliamo il salto?»

«Affermativo, ho annullato il conto alla rovescia del salto. Siamo stati
fortunati, entrambe le navi sono nel mio raggio di controllo, anche se
l'orbita della fregata si sposterà oltre il mio raggio tra dieci minuti.
Al momento, i kristang sono spaventati da una nave thuranin, una nave
madre, apparsa nel cielo senza preavviso. Li sto confondendo con
comunicazioni disturbate dall'Olandese, che non li fermeranno a lungo.
Entrambe le navi kristang si stanno preparando a saltare via con scarso
preavviso. Non si rendono conto che controllo io i loro computer.»

«Hai detto che ci sono satelliti? Che tipo di satelliti?»

«Ogni satellite è lungo cinquantotto metri, alimentato da un reattore a
fusione in grado di generare ottocentoventi megawatt di potenza maser.»

«È un sacco di roba?»

«Le vostre portaerei nucleari americane di solito generano circa
duecento megawatt, dai reattori a fissione.»

«Porca vacca.» Pensavo che i satelliti fossero cose di poco conto. La
mia idea di satellite era una scatola con pannelli solari che mi
permetteva di guardare le partite di calcio.

«Oh, molto male! Due dei satelliti si stanno preparando a sparare su una
città chiamata Mumbai, dove c'è una grande protesta contro i kristang.
Colonnello», dal tono della voce capii che Skippy era serio per un
momento, «quei satelliti e la fregata sono stati impegnati a reprimere
tentate ribellioni in tutto il pianeta. Si stimano milioni di vittime
umane. La fregata ha causato ingenti danni alle principali città con
colpi di cannone a rotaia.»

Sbattei con rabbia il pugno sul grande pulsante rosso, non pensando a
quello che stavo facendo. «Quanti kristang ci sono laggiù?»

«Millequattrocentoventitré kristang, per lo più in sette aree. Dodici
delle loro navi sono al momento in volo.»

Mi chinai in avanti con tensione, guardando lo schermo, senza rendermi
conto che il mio pugno destro era appoggiato sul grande pulsante rosso.
«Dobbiamo tenere quella nave da trasporto truppe, ma quella dannata
fregata mi piacerebbe farla saltare nel Sole. Puoi reimpostare
l'obiettivo di quei satelliti per fare in modo che colpiscano le aree
kristang e le loro navicelle?»

Ci furono bagliori luminosi in tutto il pianeta e potei vedere la nave
da trasporto truppe tremare, mentre sembrava che i corpi delle lucertole
fossero soffiati fuori dalle camere d'equilibrio. Quello stesso velivolo
sbocciò una scarica di missili, che curvarono in discesa verso il
pianeta e sfrecciarono in avanti, trasformandosi in un lampo in
striature ardenti attraverso l'atmosfera.

«Fatto», annunciò Skippy. «La popolazione kristang è ora a quota
settecentoventi. Ah! Missili che impattano e\ldots{} ok, la popolazione
è a quota centosettantadue. Aspetta. Ah! Ha preso quei bastardi!» Ci fu
un altro bagliore luminoso da un satellite. «Il satellite ha
oltrepassato l'orizzonte, ho dovuto diffrangere il raggio maser.
Centosessantaquattro.»

«Cosa?»

«Hai detto al sistema di sparare.»

«Non l'ho fatto», mi resi conto che la mia mano era appoggiata sul
pulsante. «Merda! Skippy, ti ho chiesto se potevi farlo, non ti ho detto
di farlo davvero!»

«Ops.»

«Ops? Skippy, questo è un gran cazzo di ops.»

«Non volevi neanche che facessi saltare quella fregata nel nucleo della
stella locale? Perché non lo posso annullare. Tra un paio di milioni di
anni, gli atomi di quella fregata emergeranno nella fotosfera, ma quel
Tombolo Dondolo non si potrà riassemblare.»

«Dobbiamo lavorare sulla nostra comunicazione.»

«L'ho notato.»

Spostai la mia attenzione sulla nave da trasporto truppe. «Cos'è
successo laggiù?»

«Quella nave era tutta infestata di lucertole. Bleah, disgustoso. Ho
dovuto disinfestarla con una decompressione esplosiva. Ce ne sono
quattro ancora vive laggiù, cose maledettamente testarde.» Skippy
sembrava frustrato. È probabile che fosse in grado di uccidere quelle
quattro lucertole, ma questo avrebbe potuto richiedere che danneggiasse
la nave, che invece volevo mantenere intatta.

Rabbrividii senza volerlo, pensando all'immenso potere di Skippy e a
quanto le cose sarebbero potute andare male se mi fosse accaduto di
nuovo di non prestare attenzione. Dovevo mettere un coperchio su quel
grande pulsante rosso, così, la prossima volta che avessi voluto usarlo,
avrei dovuto prima levare la copertura. «Hai il controllo completo? Le
lucertole non possono fare del male agli umani in questo momento?»

«In questo momento, immagino che siano troppo occupate a pisciarsi
addosso per pensare di fare qualcos'altro. Mi sono anche preso la
libertà di chiudere i progetti kristang di modifica dei vostri migliori
terreni agricoli a loro uso e consumo. E ho fritto tutti i loro sistemi
informatici. Qualunque cosa faranno, la faranno senza buona parte
dell'elettronica.»

«Grande, grazie.» Mi strofinai la faccia e chiusi gli occhi per un
momento. Oltre duemila lucertole, fatte fuori in un attimo, con le loro
stesse armi. Quando alzai gli occhi, quelli di Simms erano spalancati.

«È successo davvero?», chiese il maggiore, stupita.

Annuii e indicai nello schermo l'immagine della nave da trasporto
truppe, ora circondata da una nuvola di corpi di lucertole congelati
all'istante. «E più di mille kristang morti sul pianeta.» Per qualche
ragione, quassù, sembrava più naturale dire ``il pianeta'' che ``la
Terra''. Era da un po' che non vedevo il posto, comunque.

Gli altri mi annuirono attraverso il vetro. «È un buon inizio», osservò
Simms. Concordai.

Forse la comune gente perbene sarebbe inorridita di fronte a tanta morte
e distruzione. Parlando a nome della nostra allegra banda di pirati, non
avevamo simpatia per le lucertole. Se quello che avevamo sentito della
situazione sulla Terra, quello che i kristang avevano fatto alle persone
e alla biosfera, era vero, allora l'intera specie delle lucertole poteva
andare a farsi fottere.

«Skippy, puoi contattare il governo degli Stati Uniti? Ho bisogno di
parlare con una persona autorevole.»

«Certo. Un momento. Vai.»

«Pronto?» Una voce femminile, che riconobbi vagamente, uscì dagli
altoparlanti. «Chi è?» Sembrava sorpresa.

«Ehm, chi è a parlare?»

«È stato lei a chiamarmi. Dove ha preso questo numero?», chiese la voce.

E mi ricordai con uno shock dove avevo già sentito quella voce. Mai così
tirata, però, così logorata dalla disperazione.

«Skippy», dissi in un severo sussurro: «Che cazzo hai fatto?».

«Volevi parlare con il governo americano, così ho chiamato il cellulare
personale criptato della vostra presidentessa.»

Scossi il pugno in direzione del suo lucido cilindro. Avevamo seriamente
bisogno di lavorare sulla nostra comunicazione. «Signora, signora
presidentessa, qui è il colonnello, voglio dire, sergente, Joe Bishop,
dell'esercito americano, in precedenza con la X divisione di fanteria.
Abbiamo preso una nave da trasporto truppe kristang, una fregata e una
nave madre thuranin. Ora siamo in orbita e abbiamo appena ucciso tutti
tranne centosessantaquattro\ldots»

«Più i quattro sulla nave da trasporto truppe», mi riprese Skippy.

«Tutti tranne centosessantotto kristang sulla Terra o attorno alla
Terra, e abbiamo il controllo della loro nave restante e dei loro
satelliti.»

Ci fu una lunga pausa, con voci che parlavano in sottofondo. «Silenzio!
Signor Bishop, vero? Il mio aiutante militare mi sta dicendo che ci sono
stati attacchi satellitari in tutto il pianeta, ma tutti contro gli
insediamenti kristang. I thuranin non erano i patroni dei kristang? Sia
gentile, mi dica che cosa sta succedendo.» Sembrava scossa.

Feci un lungo respiro. «È una storia molto lunga, signora. I fatti
importanti sono che i kristang non hanno più il controllo della Terra e
non arriveranno rinforzi per loro, perché abbiamo chiuso il wormhole
locale.»

«Ehi! Non dimenticarti di me!», disse Skippy. «Sono l'eroe di questa
storia, tu sei solo l'impavida spalla che fornisce intrattenimento
comico.»

«Chi è che parla?», chiese la presidentessa.

Oh, che cazzo. Ero stanco. «Signora, è una lattina di birra cromata di
nome Skippy.» Feci una pausa e considerai che di sicuro era la prima
volta che qualcuno diceva una cosa del genere a un presidente americano.
«È un'intelligenza artificiale vecchia di diversi milioni di anni, che
in confronto fa sembrare i thuranin intelligenti come la melma dello
stagno.»

«Prrr.» Skippy fece una pernacchia. «Per favore, non sono così
intelligenti.»

Toccò alla presidentessa fare un lungo respiro nel telefono: «Mi sta
tornando il mal di testa».

«Come le ho detto, signora, è una lunga storia. Dobbiamo parlare.»

«Inoltre, abbiamo bisogno di pizza!», aggiunse Skippy. «Non per me, ma
la nostra allegra banda di pirati qui ha bisogno di pizza. E birre
fredde! Si festeggia, ragazzi!»

«Allegra banda di pirati?», chiese la presidentessa.

Stava venendo il mal di testa anche a me.

\hfill\break

Dopo che la presidentessa mi ebbe passato un assistente per discutere i
dettagli e dopo avere concordato i termini per portare a casa la nostra
ciurma di pirati, controllai di nuovo lo schermo principale. «Skippy,
siamo al sicuro? Ci sono navi kristang o thuranin dirette verso di noi?»
Non riuscivo ancora a interpretare gli schermi. Sembravano vuoti, non
fosse per la nave da trasporto truppe kristang e l'Olandese. E una
grande macchia, che dedussi fosse la luna.

«Non secondo le banche dati della loro ammiraglia laggiù. Era previsto
che la prossima nave madre thuranin passasse attraverso il wormhole tra
dieci giorni. Questo non accadrà. Ma è possibile che ci siano navi là
fuori e che il loro comando non pensasse che i kristang qui avessero
bisogno di saperlo.»

«Colonnello Chang, maggiore Simms?», feci il saluto a entrambi. Ora che
eravamo a casa, il mio rango teatrale era stato annullato ed ero
ritornato al mio grado regolare di sergente dell'esercito. Lo sapevamo
tutti e tre. Per quanto mi riguardava, non riuscivo a decidere se essere
deluso o sollevato. «Qualcuno deve scendere sulla Terra e informare i
nostri superiori. Skippy, puoi pilotare una navicella da remoto? Dovremo
portare la gente a Pechino, Dehli, Londra e Washington, e nessuno di noi
sa come pilotare una navicella kristang o thuranin. Il nostro Dodo è
guasto.»

«E Parigi», mi ricordò Giraud.

«E Parigi.» Mi ero dimenticato di lui, il nostro unico pirata francese.
«E dobbiamo prestare attenzione supplementare ai feriti.» Tutti quelli
che erano stati feriti durante la cattura della Fiore e l'assalto alla
base sull'asteroide avevano lasciato l'infermeria ed erano in piedi,
alcuni di loro con dispositivi medici portatili thuranin ancora
attaccati. L'avambraccio di Walorski sarebbe stato liberato dalla manica
di guarigione nel giro di una settimana. Le costole di Chang erano
tecnicamente guarite, Skippy ci aveva informato che le ossa si erano
ricomposte, anche se il tenente colonnello mi diceva di sentire ancora
un male cane quando faceva un respiro profondo. Per me, era un miracolo
della tecnologia medica thuranin, o un miracolo di Skippy, il fatto che
tutti quelli che erano stati feriti si sarebbero ripresi al 100\%.

«Posso pilotare più navicelle da remoto», rispose Skippy, come se fosse
la cosa più semplice dell'universo, «non c'è bisogno che le persone
volino su tutto il pianeta con una sola nave. Ma dopo che i mezzi
saranno atterrati e la gente sarà uscita, chiuderò le porte a chiave.
Non vorrei che voi scimmie, sapete, ci giocherellaste.»

«Capito. A parte me, c'è qualcun altro che dovremmo lasciare a bordo?»
Guardai l'allegra banda di pirati attorno al sottoscritto che stava per
disperdersi.

«Perché rimani a bordo?», chiese Chang prima che potesse farlo Skippy.

«Signore, se una nave kristang salta in orbita, io e Skippy dobbiamo
essere qui a difendere la Terra.»

«No. Voglio vedere la Terra da vicino e non solo dall'orbita!»,
insistette Skippy. «Posso controllare questa nave da qualsiasi punto
laggiù.»

Simms prese la parola: «Il pulsante di controllo delle armi è qui su
questa nave, signore». Mi ci volle un momento per capire che ``signore''
era rivolto a Skippy, non a me.

«Ah, fanculo!», la derise Skippy. «Colonnello Joe, ho appena caricato
un'app ``Grande Pulsante Rosso'' sul vostro zPhone. Se lo premete,
ovunque sulla Terra, saremo armati.»

Tirai fuori dalla tasca il telefono. «Sono di nuovo sergente, non più
colonnello. E non lo vedo.»

Skippy sospirò. «Lo chiamano smartphone perché è più intelligente
dell'utente? È nell'ultima schermata delle applicazioni, tra le due
versioni di solitario e quel grande gioco degli uccelli cui non giochi
più.»

«Sono stato un po' occupato.» Vidi l'applicazione. Era difficile che ti
sfuggisse. «Non dovrebbe essere nella prima schermata?»

«Joe, Joe, Joe. Non vogliamo che tu selezioni le armi digitando a cazzo
e distrugga per sbaglio, che so, il Canada.»

«Merda!» Tenni lo zPhone lontano da me. «Potrebbe succedere?»

«Improbabile, visto che sono io a programmare le armi, tu premi solo il
pulsante per autorizzarne l'uso. Però abbiamo avuto problemi di
comunicazione, come dici tu, quindi\ldots»

«Sergente Bishop, mi sentirei più a mio agio se tu, e, ehm», Simms si
sforzava di trovare un modo per evitare di dire il nome di Skippy, «il
signor Skippy veniste con me, per informare la nostra leadership,
compresa la presidentessa, che già conosce a quanto pare.» Non ne
sembrava felice.

Chang si chinò verso Simms. «Non sono d'accordo che il dispositivo»,
indicò Skippy, «debba essere di fatto di proprietà dell'America. In
quanto ufficiale Unef in comando\ldots»

«Non sei tu il mio ufficiale in comando, Changy-boy. Il capitano di
questo equipaggio pirata è ancora il colonnello Joe», puntualizzò Skippy
con un tono decisamente scortese. «E chiamarmi ``il dispositivo'' non è
un buon modo per entrare nelle mie grazie. Prima che tu dica
qualcos'altro di stupido, non sono di proprietà di nessuno ed è
probabile che mi verrà chiesto di visitare la Cina mentre sono qui,
perciò i tuoi scienziati avranno l'opportunità di farmi le stesse
stupide domande che mi faranno gli americani. Ora, chiedi scusa a me e
al colonnello Joe o la tua navicella potrebbe finire nel deserto del
Gobi, per errore.»

«Non c'è bisogno di scuse, signore», mi affrettai a dire. Chang, era,
dopotutto, un vero tenente colonnello e un bravo ragazzo. Cazzo.
Possibile che chiudere il wormhole e sconfiggere i kristang fosse stata
la parte più facile?

Chang fece un lieve inchino: «Mi scuso, signor Skippy, senza offesa. Era
mia intenzione assicurarmi che i diritti della Cina fossero presi nella
giusta considerazione».

Skippy sospirò di nuovo. Stava diventando un'abitudine: «Credimi, non ho
alcun interesse ad aiutare un gruppo di scimmie laggiù a guadagnare un
vantaggio sulle altre. Potete picchiarvi a vicenda con i bastoni quando
non ci sono. Se avete più cervello di un'ameba, forse vorrete
concentrare la vostra energia sulla riparazione del danno che i kristang
hanno causato al vostro pianeta. Faccio per dire».

«Signore?» Il sergente Adams rivolse la sua domanda a Chang: «Possiamo
contattare le nostre famiglie?».

Skippy parlò in privato nel mio auricolare: «Colonnello Joe, la sua
famiglia è sana e salva, ma lo stesso non vale per le famiglie e gli
amici dell'intero equipaggio. Starei attento alle comunicazioni per ora.
Ho appena detto la stessa cosa a Chang. Solo che gli ho anche detto che
il fratello di suo padre è stato ucciso dai kristang».

Attraverso il vetro, vidi Chang toccare il suo auricolare, annuire,
aggrottare le sopracciglia, poi guardare me. Si rivolse a tutta la
compagnia: «Non conosciamo ancora la situazione laggiù e, ora che siamo
tornati, siamo sotto l'autorità dei nostri governi nazionali. È
probabile che vorranno mantenere la sicurezza operativa per l'immediato
futuro. Sono sicuro che saremo tutti informati dopo l'atterraggio».

\hfill\break

La mia breve fantasia di far drammaticamente atterrare una navicella
thuranin sul prato della Casa Bianca fu infranta in fretta, poiché la
Casa Bianca e gran parte di Washington erano state danneggiate dalle
azioni dei kristang contro i traditori. Il quartier generale del governo
federale degli Stati Uniti si trovava ora a Colorado Springs, sul sito
dell'Air Force Academy, dietro fortificazioni pesanti. Skippy era per
sfrecciare sopra l'accademia a velocità ipersonica, utilizzando tutte le
capacità stealth di una navicella thuranin, per poi curvare e atterrare
dove cazzo gli pareva, per far sapere fin dall'inizio al governo
americano quanto fosse impotente. Feci notare che io ero arruolato
nell'esercito americano e avevo giurato di proteggere e difendere la
Costituzione degli Stati Uniti, perciò il governo della nazione, eletto
a norma di legge, aveva tecnicamente autorità su di me; di conseguenza,
la richiesta della presidentessa di seguire il percorso di volo
designato dall'Aeronautica militare degli Stati Uniti era, di fatto, un
ordine cui ero costretto a obbedire. Fu una conversazione lunga e
contorta, che rivelò più cose su Skippy che su di me o sugli Stati
Uniti. Skippy programmò a malincuore il pilota automatico per un profilo
d'ingresso del tutto inutilmente basso, scendendo nell'atmosfera sopra
le isole Marshall del Pacifico e accettando una ``scorta'' da una coppia
di F-22 Raptor sulla costa californiana. Skippy aveva ragione: anche se
i Raptor stavano forzando i loro limiti nella crociera supersonica,
sembrava che stessimo a malapena strisciando nel cielo. Ebbi, però, una
bella visuale del parco nazionale di Yosemite, quindi il tempo
supplementare di volo non fu del tutto sprecato.

A terra, fummo accolti da agenti dei servizi segreti che brandivano armi
automatiche, ci ordinarono di scendere piano lungo la rampa, uno per
uno. Questo fece irritare il maggiore Simms, che disse agli agenti di
abbassare le armi, prima che la super-intelligenza artificiale che aveva
spazzato via i kristang in meno di dieci secondi si arrabbiasse. Di
fronte alla loro esitazione, Simms chiese se davvero pensassero che le
munizioni delle 9 mm fossero una minaccia per Skippy. Caspita, era
incazzata. Dopo di che, le canne dei fucili furono tenute a mezz'asta.
Scesi dalla rampa, dopo gli ufficiali e il sergente scelto Adams,
essendo tornato al mio grado di sergente. Invece di infilare Skippy in
uno zaino, lo tenni nella mano sinistra davanti a me, in quella che,
nelle mie intenzioni, era una posizione d'onore. Era maledettamente
pesante, mi chiedevo se stesse regolando la sua massa per fregarmi. I
servizi segreti, grazie al loro addestramento alla scorta di sicurezza,
non avrebbero permesso a un oggetto in potenza pericoloso di avvicinarsi
alla presidentessa: fu allora che Skippy diede un taglio a quelle
puttanate e chiamò la stessa direttamente sul suo cellulare criptato.
Lei uscì dalla porta dell'ex residenza del comandante dell'accademia, e
scosse la testa in direzione dell'agente al comando, che ci fece passare
tutti, ma era evidente che andava contro il suo buonsenso. Quando gli
agenti cercarono di limitare la festa di benvenuto per me e gli
ufficiali, ci misi subito un freno con un gesto arrabbiato e feci cenno
ai soldati di proseguire. Il che strappò una risatina alla
presidentessa, che aveva fatto la stessa identica cosa nello stesso
identico istante. Da quel momento, io e la signora saremmo andati
piuttosto d'accordo, il che era una fortuna per me. Dietro di noi, i
medici si precipitarono a bordo della navicella per aiutare i membri
feriti dell'equipaggio a scendere dalla rampa.

La presidentessa salutò il maggiore Simms e altri ufficiali, ma puntò
dritto a me; la sua presenza mi spianò la strada senza dover chiedere
alla gente di spostarsi. La brezza le soffiò i capelli sugli occhi e lei
li scostò con un gesto accorto. «Signor Bishop.» I suoi occhi erano
attratti dall'insolito stemma sulla mia spalla destra. «È un paisley con
una benda sull'occhio?»

Cazzo. Ci eravamo del tutto dimenticati dell'idea di Skippy per la
nostra bandiera pirata. Ora mi sentivo un perfetto idiota. Questa donna
era responsabile dell'arsenale nucleare dell'America. Sempre che
avessimo ancora un arsenale nucleare. Cioè, le cose potevano essere
cambiate. «Una specie, signora.»

«Non è un paisley, è un paramecio», disse Skippy.

«Sto parlando con l'essere che il signor Bishop chiama Skippy?», chiese
la presidentessa, gettando di scatto uno sguardo ai suoi consiglieri che
giravano intorno con discrezione.

«L'unico e inimitabile, sono io.»

«Allora Skippy è un soprannome.» La presidentessa gli sorrise, poi mi
guardò. «Mi dica, signor Bishop, come pensa di chiamarmi?»

Ebbi un momento di panico: «Signora comandante in capo, signora?».

Rise: «Andrebbe bene, ma è piuttosto lungo. Signora presidentessa farà
al caso nostro per ora. C'è una storia dietro il paramecio?».

Come facevo a dire alla presidentessa degli Stati Uniti che una birra
cromata poteva pensare che tutta la nostra specie, compresa lei, fosse
solo un po' più intelligente dei batteri? Sarebbe stato meglio
contestualizzare. «Signora\ldots»

Skippy m'interruppe: «È una specie di scherzo. Il colonnello Joe ha
suggerito che una delle sue idee un po' meno stupide stava a significare
che la vostra specie è un po' più intelligente dei batteri. Ho accettato
a malincuore che potesse essere paragonabile a un paramecio. Su questo,
la giuria è ancora fuori per deliberare».

Oppure avrei potuto spifferare tutto come un bambino di quattro anni,
come fece Skippy! La presidentessa la prese bene. Immagino che anni di
campagne politiche le avessero fornito una dura scorza: «Certo, spero
che le faremo una buona impressione, durante la sua permanenza qui,
affinché possa paragonarci a un organismo superiore, forse alle alghe?».

«Buona fortuna con quelle», prese in giro Skippy, ma poi aggiunse
tranquillo: «Mi piace, colonnello Joe».

«Grazie», alla presidentessa riuscì un rapido sorriso. «Che significa?
Come mai la chiama ``colonnello''?», chiese rivolta a me.

«Ah, signora, era una promozione sul campo, quello che l'esercito chiama
rango teatrale. Era solo temporanea.» Guardai a disagio i galloni sulla
mia uniforme da sergente. «Ho riadeguato la mia uniforme al mio grado di
sergente dell'esercito, signora.» Le guance mi bruciavano
dall'imbarazzo, e poi ebbi un pensiero infelice. Tecnicamente, ero
tornato specialista ora? Anche la mia nomina a sergente era stata una
promozione sul campo, da parte di un esercito che poteva non esistere
più. Non avevo idea di cosa dicessero i regolamenti rispetto a questa
situazione. Cazzo, avrei dovuto chiedere a Skippy, sono certo che lui
sapesse tutto sui regolamenti dell'esercito americano. Diavolo, se
avessero voluto ritirare i miei galloni da sergente me lo avrebbero
detto. Mi piacevano i miei galloni. Me li ero guadagnati.

«Mmh», la presidentessa guardò in modo significativo il maggiore Simms,
che sembrava addolorata. «Suppongo che ci sia una lunga storia anche
dietro questo. Entrate, per favore, abbiamo del cibo per voi.» Un gemito
indisciplinato e involontario si levò dall'equipaggio pirata, forse me
compreso, e mi venne l'acquolina in bocca.

\hfill\break

Ci fu un altro intoppo mentre facevamo la fila per entrare nella stanza
che era stata allestita con un buffet. Una colazione a buffet. Stando
all'orario sull'Olandese, era pomeriggio, ma al mio stomaco non
importava. Vero cibo, cibo umano. Tutto il mio sangue doveva essere
stato reindirizzato proprio al mio stomaco, invece che al mio cervello,
perché tirai fuori con gesto distratto lo zPhone per controllare l'ora,
cosa che facevo ogni mattina a colazione. Non c'era nessuna applicazione
meteo sul mio zPhone, ovviamente.

Un agente dei servizi segreti vide la cosa nella mia mano, la mano che
non teneva Skippy, e s'irrigidì. «Signore, che cos'è?»

«Questo?», agitai lo zPhone davanti a lui. «È il mio zPhone\ldots{} oh,
ehm, l'esercito lo chiama radio tattica. I kristang ne hanno fornito uno
a ogni soldato.»

«Kristang? Questo è un dispositivo kristang?», chiese l'agente
allarmato. «Non si può portare un dispositivo di comunicazione kristang
qui dentro, devo requisirlo.»

Mi strinsi lo zPhone al petto, in una presa mortale, come se fosse un
pallone che avevo appena intercettato. «Non posso farlo, signore. Non
sto cercando di complicarle il lavoro. Dovete capire che è con questo
che posso esercitare i miei controlli da remoto, e attraverso Skippy
qui, su tutte le armi della nave madre thuranin in orbita.» Indicai con
gesto vago il cielo.

«Signore?» Era chiaro che l'agente non stava mangiando la foglia.

«Quello che ha detto», si offrì Skippy invano. «Non mente.»

Adesso era ancora più chiaro che l'agente non stava mangiando la foglia:
«Signore, avrò bisogno di esaminare quel dispositivo».

Niente da fare. Avevo visto il grande pulsante rosso friggere siti
kristang su tutta la Terra. Non esisteva che qualcuno mettesse le sue
zampe sporche sul mio telefono. Cazzo. Prima che i ruhar ci assalissero
-- ora sembrava essere passata un'eternità -- ero impaziente di
aggiornare il mio vecchio cellulare. Ora non potevo mai perdere di vista
il mio zPhone. Se avessi trovato una borsa di plastica adatta, me lo
sarei portato pure nella doccia. E no, non per guardare dei porno, come
penserebbero gli sputasentenze. «Senta, agente, ehm, signore, i kristang
non possono più usare questa cosa. È sicuro. Giusto, Skippy?»

«Al 100\% al sicuro dalle lucertole, ho cancellato dal telefono tutto il
loro patetico codice.»

Fu la presidentessa a calmare le acque: «Agente Thomas», intervenne, «il
signor Bishop può tenere il suo telefono».

«Signora presidentessa, abbiamo delle procedure da seguire per una
rag\ldots»

«Questi due», la presidentessa fece oscillare un dito tra me e Skippy,
«hanno distrutto i kristang in pochi secondi.»

«Ci avremmo messo anche meno, se non ci fosse stato quel cazzo di
satellite oltre l'orizzonte. Stupidi meccanismi orbitali», brontolò
Skippy.

«Voglio dire che, se avessero intenzione di vaporizzare tutta questa
zona, un telefono sarebbe l'ultimo dei nostri problemi», disse con
gentilezza la presidentessa e mi fece cenno di proseguire. «Davvero il
suo telefono controlla le armi su quell'astronave?»

«Ehm, sì, signora. È una lunga storia.»

«Non vedo l'ora di sentirla tutta più tardi. Nel frattempo, ho un paese
da iniziare a rimettere insieme.» Mi tese la mano e io la strinsi
stordito. Ero in qualche modo passato dalla vita in una fattoria del
Maine a una stretta di mano con la presidentessa degli Stati Uniti
d'America. Lei guardò il mio zPhone da vicino. «Forse sarebbe una buona
idea farla accompagnare da una squadra di sicurezza, per essere sicuri
che non perda quel telefono.» Mi guardò negli occhi: «Dico sul serio,
sergente Bishop».

«No, signora», balbettai. «E, ehm, avrò bisogno di trovare un
caricabatterie», dissi con voce flebile.

«Non ce n'è bisogno», intervenne Skippy, «lo sto tenendo carico io.»

«Come?», chiesi. Non aveva mai detto niente al riguardo prima.

«Unicorni e polvere magica.»

«Vaffanculo, Skippy», risposi, prima di ricordare con orrore dove mi
trovassi. «Scusi, signora presidentessa.»

Quella donna sembrava davvero divertita. «Grazie, sergente. Credo che
oggi sia la prima volta che sorrido per davvero nell'ultimo anno.»

«Sì.» In quel momento mi ripromisi di tenere la bocca chiusa ogni volta
che potevo, in futuro. «Grazie a lei, non vogliamo rubarle tempo,
signora.»

«Non vedo l'ora di parlare con voi più tardi. Nel frattempo, godetevi la
vostra colazione.» Si girò per andare, ma si voltò a guardarmi al di
sopra della spalla. «Ah, sergente?»

«Sì, signora presidentessa?»

«Grazie per avere salvato il mondo.»

\hfill\break

La colazione era buona, strabuona. Il maggiore Simms ammonì noi ex
pirati di non ingozzarci, dal momento che i nostri stomaci non erano
abituati a mangiare cibo vero. Io mi accontentai di porridge di mais, in
cui misi zucchero di canna e panna, più toast imburrato e uno, no, due
-- ok, lo confesso --, tre strisce di pancetta. Oh, cavolo era tutto
taaanto buono. E una tazza di sacrosanto, autentico caffè. Nero, caldo,
in una tazza con il logo dell'esercito. Mi veniva da piangere.

Due cose gettarono una minuscola ombra sulla mia colazione. La prima fu
il sergente dell'Aeronautica che si presentò subito dopo che mi ero
seduto. Era alta solo un metro e sessanta, aveva i capelli castani
raccolti in una coda di cavallo corta, ma il modo in cui stava in piedi,
il modo in cui parlava e l'espressione che aveva in viso mi dicevano che
era tutta d'un pezzo. Quello, e la sua pistola. Ah, e i quattro aviatori
dietro di lei, armati di M4. «Sergente Bishop? Sono il sergente scelto
Kendall, mi hanno assegnata a te, per accompagnarti. Ovunque.»

«Ah», riuscii a dire con la bocca piena di pane tostato, «siete la
squadra di sicurezza assegnata dalla presidentessa?»

«Non so nulla della presidentessa, sergente, ma il capo di stato
maggiore dell'Aeronautica mi ha dato ordini personalmente.» Sollevò il
sopracciglio per sottolineare che non succedeva tutti i giorni.
«Dobbiamo garantire la tua sicurezza, quella del tuo telefono», disse in
tono interrogativo, «e dello, ehm\ldots», indicò Skippy, che avevo
appoggiato con premura in una ciotola di ceramica sul tavolo di fronte a
me, «dello Skippy?»

«Eccomi, in tutta la mia magnificenza, l'incredibile Skippy», si inserì
la mia lattina di birra super-intelligente. «Quindi, sei la baby sitter
del colonnello Joe? Assicurati che vada a letto presto, diventa
irritabile se non fa un bel pisolino. E non fargli bere troppi succhi di
frutta.»

Lei non sorrise, nemmeno un po'. A quanto pareva, l'Aeronautica non
aveva ritenuto opportuno dotarla di senso dell'umorismo. Ottimo. Mi
avrebbe seguito ovunque e nemmeno Skippy sarebbe riuscito a farla
sorridere.

La seconda ombra sul mio banchetto era rappresentata dai tre tipi seduti
di fronte a me, che sembravano non aspettare altro che finissi di
mangiare pane tostato. Due erano dei servizi segreti dell'esercito,
l'altro doveva far parte di una delle agenzie a tre lettere, e non si
trattava dell'Epa, non so se mi spiego²⁹.

Volevano entrare subito nel vivo dell'interrogatorio, mentre assaporavo
la colazione. Il tizio della Cia riuscì a farmi saltare la mosca al naso
all'istante: «Sergente Bishop, tutte le informazioni che ha riguardo
agli eventi fuori dai confini del pianeta sono considerate top secret, e
le è proibito\ldots».

«Oh, fanculo», sparò Skippy in modo sprezzante. «Ehi, voi scimmie siete
tutte stupide per me, ma voi lo sembrate in modo particolare. Ho una
notizia: siamo arrivati qui su una nave madre thuranin che usa una
fregata kristang come ruota di scorta. Nessuno dei vostri merdosi
presunti segreti quaggiù conta più un cazzo, e non vi va proprio giù,
vero? Le uniche informazioni che vale la pena tenere segrete sono nella
testa del colonnello Joe e nelle mie banche dati. Se volete qualche
informazione da me, potete rendervi utili e rifornire la tazza di caffè
di Joe.»

Il tizio della Cia pensava che Skippy stesse scherzando. Il che fece
incazzare la mia intelligenza artificiale. «Ehi, testa di cazzo, se voi
stupide scimmie glabre volete informazioni da me, alza subito il tuo
culo pigro da quella sedia e porta una tazza di caffè al mio amico. E
vedi di darti una mossa!»

Io porsi la mia tazza di caffè. I due ragazzi dell'esercito stavano per
piegarsi in due dalle risate, quando il tizio della Cia, con la faccia
paonazza, mi portò una tazza di caffè appena fatto. Mangiai con calma
l'ultimo pezzo di pane tostato e mi rivolsi solo ai ragazzi
dell'esercito: «Signori, cosa volete sapere prima?».

\hfill\break

Dopo la colazione, che non durò per niente a lungo quanto avrei voluto,
Skippy aveva un appuntamento con un gruppo di scienziati, che aveva a
malincuore accettato per cavarsi il dente. Dal tono della sua voce,
capii che non era ansioso di essere interrogato dall'equivalente di una
manciata di batteri, anche se alcuni di loro erano batteri vincitori del
premio Nobel. «Oh, sarà una grande rottura di palle», brontolò rivolto a
me, mentre camminavamo lungo il corridoio diretti alla sala conferenze.
«Tu lo sai di essere troppo stupido per capire qualcosa dell'universo,
perciò non fai domande stupide. Queste scimmie pompose mi faranno solo
perdere del tempo. Meeerda. Meglio cavarsi il dente.»

Era chiaro che Skippy aveva passato troppo tempo con il sottoscritto.

Gli diedi ragione quando arrivammo alla sala conferenze, affollata di
scienziati, addetti alla sicurezza e apparecchiature audio-video.
Avevano preparato un tavolo, circondato da microfoni e telecamere, su
cui posai Skippy con attenzione, augurandogli buona fortuna.

Poi arrivò il momento del mio interrogatorio, lungo il corridoio,
scortato dal sergente scelto Kendall. Si sarebbe svolto in una stanza
più piccola, senza sedie di lusso, e l'esercito, l'Aeronautica, gli
ufficiali della Marina e la gente della Cia non erano in vena di
scherzare. Il tizio della Cia continuava a fissarmi, era ancora
incazzato per quello che Skippy aveva detto a colazione. Almeno il caffè
era caldo e appena fatto. Ogni volta che ne avevo bevuto un sorso, lo
avevo fissato dritto negli occhi. Se gli sguardi potessero uccidere,
sarei morto. Per fortuna, uno di noi due aveva distrutto un intero
gruppo kristang, e -- indovinate un po'? -- non era lui.

Iniziai raccontando per sommi capi la mia storia, dal giorno in cui
eravamo partiti da Campo Alfa. Poi mi misero intensamente sotto torchio
per conoscere i dettagli, e non solo riguardo Skippy. Fornii loro
informazioni strabilianti sui kristang, i wormhole, gli Anziani, tutto
quello che avevo imparato dalla Bürgermeister. Poi, vollero conoscere la
situazione su Paradiso, la disposizione delle nostre forze in tutto il
pianeta, lo status delle forze ruhar, kristang e thuranin. Tutte cose di
cui sapevo poco, purtroppo. Questo non impedì loro di farmi le stesse
domande più e più volte, finché non cominciai ad avvertire il mio
didietro andare a fuoco. Ebbi la sensazione di deluderli come ufficiale
militare: avrei dovuto fare di più per raccogliere informazioni. Ma come
avrei potuto farlo, dopo essere stato prima prigioniero dei kristang,
poi dei ruhar ed essere sgattaiolato in lungo e in largo per evitare
tutti loro? Non ne avevo idea. Quello che sapevo era cosa si sarebbe
aspettato l'esercito da un colonnello, e io non stavo rispettando quegli
standard.

A parte, sapete, salvare il mondo.

Avrei fatto conto su quello.

\hfill\break

Durante una pausa per andare in bagno, mi stavo lavando le mani accanto
a uno degli ufficiali dell'esercito, un certo colonnello Landry che fino
a quel momento mi aveva trattato abbastanza bene. Lo guardai nello
specchio e gli chiesi: «Signore, le ho detto che abbiamo avuto i
biscotti della fortuna, abbiamo sentito cosa stavano facendo i kristang
ai nostri terreni agricoli e ai Grandi Laghi, ma non ho sentito molti
dettagli al mio livello. Perché il governo federale è qui a Colorado
Springs?».

«Hai sentito che ci sono state proteste in tutto il mondo, quando la
gente ha saputo cosa ci chiedevano i kristang?» Landry fece una pausa e
io annuii. «Quando le proteste a Washington si sono fatte serie e hanno
iniziato ad attirare persone da tutto il paese, abbiamo portato elementi
del XXVIII fanteria per controllare la folla, speravamo di contenere le
proteste abbastanza da impedire che i kristang si preoccupassero e
decidessero di reagire. La necessità di controllare la folla ha
comportato l'utilizzo di gas lacrimogeni, proiettili di gomma, cannoni
ad acqua, tutte cose per cui il XXVIII non è addestrato. Penso che
l'idea fosse sfruttare i loro Stryker per intimidire la folla, ma hanno
ottenuto l'effetto opposto: i manifestanti hanno iniziato a lanciare
bottiglie e molotov contro le nostre truppe. Abbiamo evacuato la
presidentessa e il congresso, prima a St. Louis, poi qui. Ci sono stati
alcuni incidenti, e la folla era del tutto fuori controllo. I kristang
ci hanno consigliato, caldamente, di sparare ai manifestanti. Quando il
XXVIII si è rifiutato di usare munizioni vive, i kristang ci hanno
ordinato, con spudoratezza, di inviare la IX forza aerea per bombardare
la folla. Quando ci siamo rifiutati di farlo, i kristang hanno colpito
dall'orbita la base aerea di Shaw, cancellando la IX dalla faccia della
Terra. Poi hanno usato una specie di arma laser a microonde, hanno
ucciso metà delle persone a Washington, compresi gli uomini del XXVIII.
Da lì le cose sono precipitate.»

Landry finì di lavarsi le mani e afferrò un tovagliolo di carta, vidi
che i suoi arti erano percorse da un lieve tremito. Fece un respiro
profondo, come se stesse cercando di decidere: «C'è ancora molto che
dobbiamo sapere da te, ma, comunque tu abbia fatto, ci hai tirati fuori
da un casino infernale, sergente, e l'esercito non lo dimenticherà.
Prima che tu apparissi in orbita e colpissi i kristang», fece un lungo
respiro, «stavamo esaminando le ipotesi per il destino della Nca»,
intendeva l'Autorità di comando nazionale, «e non c'erano buone
prospettive. Il meglio in cui speravamo era una specie di sopravvivenza,
per la razza umana. Ci hai tirato fuori da una situazione senza via
d'uscita. Non lo dimenticare».

\hfill\break

Tornai all'interrogatorio, cercando di rispondere alle domande nel modo
più completo possibile. Dopo molte ore, erano diventate ripetitive, e
ormai replicavo quasi in automatico, quando il tizio della Cia
s'intrufolò a sorpresa.

«Sergente, possiamo isolare il dispositivo?»

«Scusi, quale dispositivo?» Merda, avrei dovuto prestare più attenzione.
A mia difesa, c'è da dire che mi avevano interrogato per nove ore piene,
con solo brevi pause bagno e uno spuntino.

«L'intelligenza artificiale, l'intelligenza artificiale aliena.»

«Intende Skippy? Non lo chiamerei dispositivo, non gli piace. Cosa vuol
dire isolare? Tipo ignorarlo? È molto insistente.»

«Isolarlo, tipo in un caveau di piombo o sotto il monte Cheyenne, o
dentro una gabbia di Faraday, qualcosa che lo isoli dall'accesso a
sistemi elettronici esterni», spiegò l'agente nel tono che useresti con
un bambino. «Sergente, capisce che questa intelligenza artificiale
rappresenta un potenziale rischio per la sicurezza?»

«Senta, lei non conosce Skippy come lo conosco io. L'ho visto deformare
lo spazio-tempo e far saltare un'astronave un terzo di anno luce fuori
rotta. Ed è quello che fa per hobby. È probabile che potreste farlo
cadere in un vulcano e questo non lo riscalderebbe neppure, né
interferirebbe con la sua capacità di friggere ogni sistema elettronico
su questo pianeta. Lo farebbe solo incazzare. Io non lo farei.»

«Dobbiamo prendere in considerazione\ldots»

Ora il tizio della Cia stava facendo incazzare me. «No, gliel'ho già
detto», mi astenni dall'aggiungere ``idiota''. «Skippy ha interrotto
solo in via temporanea il funzionamento del wormhole locale, che si
riavvierà o si resetterà, o qualsiasi altra cosa, in meno di un mese. Se
lasciamo che accada, ci saranno delle lucertole arrabbiate che
l'attraverseranno e si chiederanno cosa diavolo stia succedendo con la
Terra. Skippy non è un rischio per la sicurezza, è la nostra unica
speranza per la sicurezza. Dobbiamo fare quello che gli ho promesso e
tornare là fuori perché possa chiudere il wormhole per sempre e sbattere
fuori tutte le altre specie. Poi lo aiuteremo a trovare questo
Collettivo, qualunque cosa sia.»

Il tizio della Cia mi guardò torvo, a quanto pare eravamo incazzati
l'uno con l'altro: «È una promessa che non aveva l'autorità di fare,
sergente, e ora che è qui\ldots».

Autorità? Dov'era quell'idiota quando avevo dovuto inventarmi il piano
in corso d'opera? Fare quella promessa a Skippy aveva liberato la Terra
dalla schiavitù delle lucertole! Ora i ragazzi dell'esercito e della
Marina stavano lanciando sguardi denigratori al tizio della Cia e
bisbigliando l'uno con l'altro, questo mi rese audace. «Siamo qui solo
perché avevamo bisogno di rifornimenti a bordo dell'Olandese Volante e
ho convinto Skippy che le nostre perdite erano troppo numerose per
continuare la missione. Si è occupato dei kristang qui perché gli
mettevano i bastoni tra le ruote. Vuole mettergli i bastoni tra le
ruote?»

Il colonnello Landry intervenne per mettere a tacere il tizio della Cia:
«Sergente, comprendiamo la situazione e ne stiamo discutendo ai massimi
livelli. Le informazioni che ci fornisci ispireranno tale decisione.
Ora, vorrei tornare a qualcosa che questo, ehm\ldots», controllò i suoi
appunti, «questa Bürgermeister ti ha detto\ldots»

\hfill\break

C'era dell'alcol alla reception dopo il mio primo giorno di tormentato
interrogatorio, una sostanza di cui non avevo goduto da prima di
lasciare la Terra sull'ascensore spaziale. Fu una pessima idea ordinare
rum e Coca. Pensavo che, dopo una lunga giornata di interrogatori, mi
servisse qualcosa con dentro della caffeina, di qui la Coca-Cola. Il rum
suonava così bene che dovevo provarci. Ed era buono. Anche pericoloso.
Dopo un delizioso sorso, lo appoggiai dietro una pianta e chiesi al
barista di darmi un Club soda e lime. Il sergente scelto Kendall annuì
in segno d'approvazione. Dopotutto, avevo in tasca il grande pulsante
rosso. Finché avevo quella responsabilità, ero in servizio 24 ore su 24
e non avrei dovuto bere niente di più forte dell'acqua.

Sembrava che la gente non sapesse cosa farsene di me, un sergente
semplice, che non avrebbe neppure dovuto essere lì, a parte la faccenda
di aver salvato il mondo. Il che faceva sentire le persone a disagio,
non sapendo cosa dirmi. Alla fine, andai al bar a parlare con il
barista, un soldato dell'esercito del Texas di nome Matt. Tra la sua
parlata strascicata da cowboy e la mia pronuncia nasale del Sud-Est,
dovevamo scandire i termini per capirci. Mi fece sentire la nostalgia di
Cornpone e Ski e del Sergente Koch e dei ragazzi del mio vecchio gruppo
di fuoco, e mi chiesi come se la passassero loro e il resto dell'Unef,
da prigionieri dei ruhar.

Mi guardai attorno nella stanza e vidi un gruppo di scienziati che
parlavano e lanciavano sguardi nella mia direzione. Uno di loro, un tipo
alto, con i capelli castani, una dolcevita, giacca di tweed e mocassini
marroni, continuava a sorridermi. Sembrava il Professore Nerd numero 5
di ogni film di fantascienza che avessi mai visto. Il tipo di persona
che riusciva a fare menzione del suo premio Nobel fin dal primo minuto
di ogni conversazione. Era intelligente, rispettato ed era evidente che
lì era nel suo elemento più di quanto non lo fossi io. Ero invidioso, e
fin dal primo istante odiai quel tizio dall'altra parte della stanza.
Poi lui e il sorrisetto compiaciuto sulla sua faccia vennero verso il
bar.

«Posso portarle qualcosa, dottor Constantine?», chiese Matt.

«No, grazie», disse Constantine, senza nemmeno un rapido sguardo nella
direzione di Matt, «voglio parlare con il sergente Bishop.»

Nella mia limitata esperienza, quando un civile si riferisce a te per
rango e il tuo grado non è almeno quello di capitano, o superiore, a
volte lo fa per sottolineare quanto tu sia in basso nella piramide. Quel
tizio disse ``sergente'' come se si trattasse di una zanzara molesta da
schiacciare sul muro. «Spero che lei sappia quanto è stato fortunato ad
aver incontrato l'Aaib.» Lo pronunciò come ``Abe'', come in Abe Lincoln.

«Incontrato cosa?»

«L'A-a-i-b. Advanced Artificial Intelligence Being, l'Essere di
intelligenza artificiale avanzata. Ci è stato ordinato di non utilizzare
il termine ``dispositivo'' per riferirci a esso.»

«Dovrebbe anche evitare con cura di usare il termine ``esso''. Skippy
preferisce ``lui'', anche se penso che sia per il nostro bene. E, dal
momento che non abbiamo nessun tipo di intelligenza artificiale qui da
noi, ``avanzata'' non sarà ridondante?»

Constantine mi gettò uno sguardo che era un misto tra il suo solito
sorrisetto e un cipiglio: «Sergente, non apprezzo il suo atteggiamento
irriverente e, a essere franco, poco professionale. Ha avuto interazioni
con un essere d'immensa potenza che erano inappropriate e pericolose.
Pericolose non solo per lei, ma per tutta l'umanità. Sarebbe meglio per
tutti coloro che sono coinvolti se non avesse più contatti con l'Aaib.
Avrebbe dovuto portare il disp\ldots», inciampò sulle sue parole,
«l'Aaib, alle autorità competenti su Paradiso, così le interazioni
avrebbero potuto essere condotte da personale qualificato.»

Tossii e bevvi un sorso del mio delizioso succo di soda e lime. Questo
mi diede abbastanza tempo per trattenermi dal prendere a pugni in faccia
il pomposo idiota. «Le autorità competenti? Su Paradiso? Che -- mi pare
che non siate aggiornati -- sono i nostri presunti nemici, i ruhar. Le
truppe umane venivano fatte prigioniere. E per ``qualificato'' presumo
intenda qualcuno come lei.»

Annuì.

«Qualificato per via della sua vasta esperienza con le intelligenze
artificiali avanzate costruite milioni di anni fa dagli alieni, che da
allora hanno trasceso la loro esistenza fisica. Ah, aspetta, non può
essere lei, non ha nessuna esperienza con esseri come quello. Chi ce
l'ha?» Mi pizzicai il mento, come se fossi perso nei pensieri. «Ci deve
pur essere qualcuno che ha un contatto stretto con questo essere
avanzato, qualcuno che ha lavorato con questo essere avanzato per
fuggire dal territorio nemico, catturare non una, ma due astronavi
nemiche, vaporizzare più di una dozzina di altre astronavi nemiche,
chiudere il wormhole locale e annientare il dominio dei kristang sulla
Terra. Chi potrebbe essere? Lei? No, ho idea che non sia lei.»

«Il suo costante atteggiamento irriverente è proprio il motivo per cui
non dovrebbe\ldots»

Una serie di porte all'estremità della stanza si aprirono e un gruppo di
ufficiali, in rappresentanza di tutti e cinque i rami in uniforme,
attraversò la stanza, dirigendosi verso di me. Per quanto mi sforzassi,
non riuscivo a immaginare cosa ci facesse un ammiraglio della guardia
costiera a Colorado Springs. La maggior parte di loro si mosse verso
l'angolo in cui c'era un tizio che riconobbi come il capo di gabinetto
della Casa Bianca, ma due ragazzi dell'Aeronautica vennero verso di me,
spingendo con cautela un carrello con uno spesso cuscino rosso e Skippy
sopra. «Sergente Bishop», iniziò a dire uno di loro.

«Colonnello Joe! Com'è andata?» La voce di Skippy sembrava stanca, cosa
che avrebbe dovuto essere impossibile. Forse dovevo scroccare un po' di
elio-3 metallico, qualsiasi cosa fosse.

«Abbastanza bene, Skippy. È bello rivederti. Com'è andato il tuo
interrogatorio?» Non stavo mentendo, mi era mancato molto il suo lato
irascibile.

«Come previsto, è stato epico, storico, livelli galattici di
scoglionamento. Oggi abbiamo davvero testato il limite estremo di quanto
te le possano far girare, fino al livello quantico. Dopo essere stato
interrogato da molti dei vostri scienziati premi Nobel, mi sto
ricredendo, non so se la vostra specie si meriti che vi chiami batteri.
Cazzo, voi scimmie siete ottuse come una stella di neutroni. Vabbè, sto
interrompendo qualcosa?»

Indicai Constantine con il pollice: «No, stavo parlando con Gongolo
McPippas qui».

«Signor Bishop, mi sta dando rag\ldots», ansò Constantine.

«Scusa, volevo dire il dottor Gongolo McPippas.»

«Ah, sì», Skippy rise, «ho avuto il sommo dispiacere di parlare con il
dottor Mcpippas prima.»

Le orecchie di Constantine erano rosso fuoco, non riuscivo a capire se
per rabbia o imbarazzo, o entrambi: «Questo è il genere di\ldots».

«Ora stai zitto», lo ammonì Skippy. «Non puoi sederti al tavolo dei
grandi e parlarmi di nuovo, finché non risolvi l'equazione che ti ho
appena mandato sul telefono. Sbrigati ora, fai il bravo ragazzo.»

Constantine tirò fuori il suo telefono con fare sospettoso, poi le
sopracciglia gli arrivarono all'attaccatura dei capelli, mi sparò uno
sguardo che avrebbe potuto sciogliere un iceberg e se ne andò mormorando
eccitato.

«Che equazione gli hai mandato?»

Skippy fece una pernacchia: «Che cazzo ne so. Ho messo insieme un botto
di super-stringhe di stronzate. Sembra bella, però».

Risi: «Skippy, a volte sai essere davvero malvagio». Cazzo. Vidi
addirittura un accenno di sorriso sulla faccia del sergente scelto
Kendall. Forse il senso dell'umorismo ce l'aveva.

«Ehi, lo terrà buono buono, con la mente impegnata per giorni.»

«Con le mani impegnate, direi.»

«Dipende da quanto si eccita, eh?»

Questa fece ridere anche il sergente scelto Kendall.